\section{Conclusion}

The two companion papers establish the surrounding classification
problem.

For the canonical bipartite double cover
\[
X=KC(L(P)),
\]
the full automorphism group is
\[
\operatorname{Aut}(X)\cong S_5\times S_3.
\]
The complementary \(S_3\) acts on a three-member orbit of regular
\(L(P)\)-quotient frames.

For the distinct native graph
\[
G_{60}\cong\operatorname{AT4val}[60,6],
\]
the earlier automorphism paper proves
\[
\operatorname{Aut}(G_{60})
\cong
S_5\times_{C_2}D_8
\]
through complete native enumeration, internal group invariants, and an
explicit \(480\)-element group isomorphism.

The present paper addresses the geometric question left behind by that
classification. It recovers the Klein-four kernel directly in the
native graph and determines the structures through which the abstract
group acts.

\subsection{The native register}

The full automorphism group contains a free normal subgroup
\[
V_4^{\mathrm{nat}}
=
\{1,a,b,ab\}.
\]

Its action partitions the sixty native vertices into fifteen chambers
of size four:
\[
60=15\cdot4.
\]

The quotient is
\[
G_{60}/V_4^{\mathrm{nat}}
\cong
L(P).
\]

Thus the Petersen line graph is not only a distant source of symmetry
or an external recognition model. It is the exact quotient of the
native graph by its hidden four-state register.

The native permutation space decomposes into four
fifteen-dimensional character sectors. The full automorphism group
fixes \(a\) and exchanges \(b\) with \(ab\), producing the intrinsic
grammar
\[
a\mid\{b,ab\}.
\]

The equality of the \(b\)- and \(ab\)-character spectra is therefore
forced by symmetry.

\subsection{The five-point register}

The fifteen chambers admit an invariant partition into five triples.
Each triple is independent in the carrier quotient
\[
L(P).
\]

Under the explicit Petersen coordinate bridge, these triples become
five three-edge matchings partitioning the fifteen Petersen edges:
\[
E(P)
=
M_0\sqcup M_1\sqcup M_2\sqcup M_3\sqcup M_4.
\]

Every native automorphism preserves this five-block system. The
induced action is the full symmetric group:
\[
\operatorname{Aut}(G_{60})
\longrightarrow
S_5.
\]

Its kernel is exactly the native Klein group. Hence
\[
\operatorname{Aut}(G_{60})/V_4^{\mathrm{nat}}
\cong
S_5.
\]

The extension splits, and the quotient action on the kernel is
controlled by sign:
\[
\operatorname{Aut}(G_{60})
\cong
V_4\rtimes_{\operatorname{sgn}}S_5.
\]

This is the native realization of the earlier fiber-product theorem
\[
S_5\times_{C_2}D_8
\cong
V_4\rtimes_{\operatorname{sgn}}S_5.
\]

The resulting decomposition is equivariant:
\[
60
=
5
\cdot
3
\cdot
4.
\]

The three factors now have exact meanings:

\[
5
=
\text{matching blocks},
\]
\[
3
=
\text{chambers per matching block},
\]
and
\[
4
=
\text{native Klein states per chamber}.
\]

\subsection{Registered closure}

The quotient voltage coordinates are gauge dependent, but cycle
holonomy is intrinsic.

The shortest nontrivial holonomy occurs on quotient pentagons. Six
carry \(b\)-holonomy and six carry \(ab\)-holonomy. Every such
pentagon lifts to two disjoint native ten-cycles.

These twenty-four cycles are exactly the registered ten-cycles. Their
quotient positions close after five steps, while their native states
close after ten.

They form one full automorphism orbit. Their stabilizers are
\[
D_{10},
\]
and the stabilizer of one answering pair is
\[
D_{10}\times C_2.
\]

Registered closure is also visible without first consulting the
quotient projection. It has the exact fourth-order signatures
\[
x_C^{\mathsf T}A^4x_C=880
\]
and
\[
x_C^{\mathsf T}L^4x_C=2160.
\]

Within the third-moment frontier, the registered class is precisely
the class excluding two external relay paths and their sixteen
associated four-step returns.

\subsection{The foreign Klein layer}

The \(120\) nonregistered frontier cycles form one automorphism orbit.
Their stabilizers are foreign Klein four-groups.

There are fifteen distinct foreign Klein groups, each stabilizing
eight frontier cycles. Distinct foreign groups intersect only in the
identity, and none is conjugate to the native Klein group.

Their chamber-incidence graph is
\[
5K_3.
\]

The five triangles are exactly the five matching blocks recovered
independently from the carrier quotient. Within each block, carrier
adjacency is absent because the corresponding Petersen edges are
pairwise disjoint. Foreign-Klein incidence supplies all three internal
relations.

This produces two complementary structures on the same fifteen
chambers:

\[
\text{carrier adjacency}
=
\text{shared Petersen endpoint},
\]
and
\[
\text{foreign incidence}
=
\text{common recovered matching}.
\]

The carrier separates each matching. The foreign Klein layer completes
its triangle.

\subsection{What has been explained}

The earlier paper determines which group acts:
\[
\operatorname{Aut}(G_{60})
\cong
S_5\times_{C_2}D_8.
\]

The present paper explains how that group is realized.

The Klein-four kernel is the native four-state chamber register.

The \(S_5\)-coordinate is the full permutation action on five invariant
Petersen matching blocks.

The shared \(C_2\)-coordinate is permutation parity.

The distinguished central direction is \(a\).

The exchanged hidden directions are \(b\) and \(ab\).

The dihedral coordinate is reflected geometrically in registered
decagon stabilizers, answering-pair exchange, and the sign-coupled
action on the native Klein kernel.

The classification is therefore no longer only abstract:
\[
\boxed{
\operatorname{Aut}(G_{60})
\cong
V_4\rtimes_{\operatorname{sgn}}S_5
}
\]
is realized by a concrete hierarchy of chambers, matchings,
characters, cycles, and stabilizers.

\subsection{Remaining questions}

Several finite structural questions remain natural.

\begin{enumerate}
    \item Determine whether the two recovered \(S_5\)-complements
    exhaust all complement conjugacy classes.

    \item Give a conceptual, non-enumerative derivation of the five
    Petersen matching partition from the quotient action.

    \item Derive the fourth-moment recognition law directly from the
    \(V_4\)-character operators and the quotient-pentagon geometry.

    \item Classify the fifteen foreign Klein groups abstractly inside
    \[
    V_4\rtimes_{\operatorname{sgn}}S_5
    \]
    without first constructing the frontier cycles.

    \item Determine whether analogous native and foreign Klein
    geometries occur in the other invariant double-cover classes over
    \(L(P)\).
\end{enumerate}

These questions concern deeper explanations and extensions of the
finite structure. They do not reopen the classification established
here.

\subsection{Final boundary}

The present work proves a theorem about one exact finite graph.

It does not derive a physical theory, identify graph moments with
energy, identify registered closure with mass, or assign a physical
meaning to permutation parity.

Its earned conclusion is:

\begin{quote}
The native graph \(G_{60}\) carries an intrinsic
\(5\times3\times4\) register. Its visible \(S_5\)-motion acts on five
Petersen matching blocks, its hidden \(V_4\)-state acts within each
chamber, and the two layers are coupled exactly by parity.
\end{quote}
